% titel: Blatt 3, A2: Reaktionsgleichungen ausgleichen
% skript: Mortimer Kap. 4.3 (Reaktionsgleichungen, Oxidationszahlen)

% Teil I: Werkzeuge, kein Lösungsweg.

Das Ausgleichen (Balancieren) einer Reaktionsgleichung stellt sicher, dass für
jedes Element rechts und links gleich viele Atome stehen (Massenerhaltung) und
bei Ionen-/Redoxreaktionen die Ladung erhalten bleibt.

\begin{werkzeugkasten}[Ausgleichen Schritt für Schritt]
\begin{itemize}
  \item Jedes Elementsymbol einzeln zählen, idealerweise in einer Tabelle
        (links vs.\ rechts).
  \item Faktoren vor die \emph{ganzen} Verbindungen setzen, nie die Indizes
        der Formel verändern.
  \item Wirksame Reihenfolge: zuerst Metalle, dann Nichtmetalle (ausser
        $\mathrm{O}$ und $\mathrm{H}$), dann $\mathrm{O}$ und zum Schluss
        $\mathrm{H}$; so bleiben die letzten beiden Elemente als Freiheitsgrade.
  \item Bei Redoxreaktionen die halbreaktionensmethode nutzen: Reduktion und
        Oxidation getrennt ausgleichen (Ladung mit $\mathrm{e^-}$ und ggf.
        $\mathrm{H_2O}$/$\mathrm{H^+}$ bzw.\ $\mathrm{OH^-}$), dann addieren.
  \item Ergebnis kontrollieren: jedes Element und die Nettoladung müssen
        beidseitig übereinstimmen.
\end{itemize}
\textit{Skript:} Mortimer Kap. 4 (Ausgleichen, Redoxreaktionen), Übung 3/2.
\end{werkzeugkasten}

\begin{achtung}[Darauf achten]
\begin{itemize}
  \item Indizes stehen lassen: $\mathrm{TiO_2}$ bleibt $\mathrm{TiO_2}$; der
        Faktor gehört vor die Formel.
  \item Ladungen bei Ionen ($\mathrm{S_2O_3^{2-}}$, $\mathrm{I^-}$) in der
        Bilanz berücksichtigen.
  \item Die einfachsten ganzzahligen Koeffizienten wählen; wenn möglich durch
        den gemeinsamen Teiler kürzen.
  \item Nicht zu früh aufhören: Die Gleichung erst als vollständig ausgeglichen
        betrachten, wenn \emph{alle} Elemente und die Ladung passen.
\end{itemize}
\end{achtung}
